\chapter{Methodology}
\label{chap:methodology}

This chapter describes the research design, engineering models,
design-space exploration, optimization method, native CAD generation, and
validation approach. The bracket is used as the complete proof of concept,
while the padeye and stabilizer demonstrate the extension of the method to
different structural forms and load paths.

\section{Research Design}
\label{sec:research-design}

The project used a design-and-evaluate research method. It created an
engineering software artifact, tested the artifact, applied it to three case
studies, and evaluated the resulting analytical and CAD evidence against
declared requirements. The work did not use experimental human subjects or a
statistical population. Its evidence came from software behavior, analytical
calculations, native CAD artifacts, and a manually configured solver report.

The work proceeded through four technical phases:

\begin{enumerate}
    \item \textbf{Engineering definition:} define parameters, bounds, units,
    baselines, materials, loads, constraints, fixed geometry, and boundary
    roles.
    \item \textbf{Search:} evaluate baselines, sample the design spaces, and
    run three bounded Nelder--Mead starts for each part.
    \item \textbf{CAD verification:} rebuild selected designs in Fusion,
    calculate physical properties, export artifacts, and check boundary roles.
    \item \textbf{Validation preparation:} package baselines and finalists,
    prepare manual Static Stress studies, and record the evidence required for
    engineering acceptance.
\end{enumerate}

The phases follow the evidence tiers in Figure~\ref{fig:tiered-evidence}. A
later tier can reject an earlier candidate. For example, CAD failure can
reject an analytically feasible design. Corrected FEA can reject a design that
passes both analytical and CAD checks.

\subsection{Workflow Architecture}
\label{sec:architecture}

The workflow contains four functional layers: a shared engineering
definition, the analytical search, the Fusion CAD handoff, and the resulting
evidence record. Figure~\ref{fig:system-architecture} shows how these layers
support the study.

\begin{figure}[H]
\centering
\begin{tikzpicture}[
    node distance=0.72cm,
    box/.style={
        rectangle,
        rounded corners,
        draw=nemoblue,
        fill=nemolight,
        thick,
        text width=0.72\textwidth,
        minimum height=1.0cm,
        align=center,
        font=\small
    },
    data/.style={
        rectangle,
        draw=nemoteal,
        fill=green!5,
        thick,
        text width=0.72\textwidth,
        minimum height=1.0cm,
        align=center,
        font=\small
    }
]
\node[data] (manifest) {\textbf{Engineering definition}\\
Parameters, units, bounds, baselines, materials, loads, constraints, and
boundary roles};
\node[box,below=of manifest] (engine) {\textbf{Analytical study}\\
Sampling, evaluation, optimization, failure handling, and finalist selection};
\node[box,below=of engine] (bridge)
{\textbf{CAD handoff}\\Versioned and correlated request/response};
\node[box,below=of bridge] (fusion)
{\textbf{Autodesk Fusion}\\Native CAD, physical properties, STEP, and
boundary metadata};
\node[data,below=of fusion] (records) {\textbf{Evidence records}\\
Analytical results, validation sets, CAD artifacts, and study review};

\draw[arrow] (manifest) -- (engine);
\draw[arrow] (engine) -- (bridge);
\draw[arrow,<->] (bridge) -- node[right,font=\scriptsize]{versioned JSON files} (fusion);
\draw[arrow] (fusion) -- (records);
\draw[arrow] (engine.west) -- ++(-0.65,0) |- (records.west);
\end{tikzpicture}
\caption{NEMO study workflow and evidence flow}
\label{fig:system-architecture}
\end{figure}

The analytical and CAD stages run in separate Python environments. Their
separation allows the search to use the required scientific dependencies while
Fusion retains its native application environment. A small versioned exchange
passes candidate parameters to Fusion and returns geometry properties and
boundary evidence. Autodesk defines the application lifecycle for Fusion
scripts and add-ins \cite{autodeskAddins}.

\section{Engineering Models}
\label{sec:engineering-models}

\subsection{Common Engineering Definition}
\label{sec:part-definition}

Each component was represented by one engineering definition containing its
design variables, units, lower and upper bounds, baseline dimensions, material
properties, design load, acceptance criteria, fixed geometry, and required
boundary roles. Most dimensions are expressed in millimetres; the stabilizer
spar locations are percentages of local chord. Lengths were converted to
metres for the SI calculations, while stress and displacement were reported in
MPa and millimetres. The same definition governed analytical evaluation,
optimization, candidate selection, and CAD generation so that a change in
load, bound, or unit was not applied differently at successive stages.

\subsection{Bracket}
\label{sec:bracket-definition}

The bracket represents a foundation for a 50~kg auxiliary unit. The design
load uses a project dynamic factor of 3:

\begin{equation}
F_d=mg\gamma_d=50(9.81)(3)=1471.5~\mathrm{N}.
\label{eq:bracket-load}
\end{equation}

The dynamic factor is an engineering envelope. It is not a modal, shock, or
vessel-motion result. The analytical calculation retains 1471.5~N. The current
Fusion input dialog does not accept the decimal force value used in this
workflow. The manual study therefore uses the conservative integer value
1472~N, a difference of approximately 0.034\%.

The project-defined Aluminum 6061-T6 properties are density
2700~kg/m$^3$, yield strength 276~MPa, Young's modulus 68.9~GPa, and
Poisson ratio 0.33. Table~\ref{tab:bracket-parameters} gives the variables.

\begin{table}[H]
\centering
\caption{Bracket design variables}
\label{tab:bracket-parameters}
\begin{tabularx}{\textwidth}{Xcrrr}
\toprule
\textbf{Variable} & \textbf{Unit} & \textbf{Lower} & \textbf{Baseline} &
\textbf{Upper}\\
\midrule
Baseplate length & mm & 100 & 150 & 200\\
Baseplate width & mm & 80 & 100 & 150\\
Baseplate thickness & mm & 4 & 8 & 15\\
Rib height & mm & 20 & 40 & 60\\
Rib thickness & mm & 3 & 6 & 10\\
Fillet radius & mm & 2 & 5 & 10\\
\bottomrule
\end{tabularx}
\end{table}

The fixed geometry includes four 10~mm mounting holes and two ribs. The
screening acceptance criteria are $FOS\geq2.5$ and
$\delta\leq0.5$~mm. These limits are project criteria and do not replace a
governing marine rule set.

\subsubsection{Analytical formulation}
\label{sec:bracket-model}

All geometric dimensions enter the equations in meters. The approximate
volume contains the rectangular baseplate, four-hole subtraction, two
triangular ribs, and a fillet contribution:

\begin{equation}
V_b =
LWt_b
-n_h\pi\left(\frac{d_h}{2}\right)^2t_b
+n_r\frac{Lh_rt_r}{2}
+n_r\frac{\pi r_f^2}{4}\max(0.35W,t_r).
\label{eq:bracket-volume}
\end{equation}

The bracket mass is

\begin{equation}
m_b=\rho V_b.
\label{eq:bracket-mass}
\end{equation}

The equivalent cantilever moment is

\begin{equation}
M_b=F_de,\qquad
e=\max(0.03,0.35L+0.02).
\label{eq:bracket-moment}
\end{equation}

The effective section modulus combines the baseplate and rib terms:

\begin{equation}
Z_{\mathrm{eff}}=
\frac{Wt_b^2}{6}
+0.65n_r\frac{t_rh_r^2}{6}.
\label{eq:bracket-section-modulus}
\end{equation}

The stress estimate includes a fillet-dependent concentration factor:

\begin{align}
K_t&=\max\left(1.08,1.35-\frac{r_f[\mathrm{mm}]}{40}\right),
\label{eq:bracket-kt}\\
\sigma_b&=\frac{M_b}{Z_{\mathrm{eff}}}K_t.
\label{eq:bracket-stress}
\end{align}

The effective second moment of area is

\begin{equation}
I_{\mathrm{eff}}=
\frac{Wt_b^3}{12}
+0.45n_r\frac{t_rh_r^3}{36},
\label{eq:bracket-inertia}
\end{equation}

and the tip-equivalent displacement is

\begin{equation}
\delta_b=\frac{F_de^3}{3EI_{\mathrm{eff}}}.
\label{eq:bracket-displacement}
\end{equation}

The yield-based factor of safety is

\begin{equation}
FOS_b=\frac{\sigma_y}{\sigma_b}.
\label{eq:bracket-fos}
\end{equation}

The equations preserve the expected broad trends. Increased plate or rib
thickness increases mass and stiffness. Increased rib height strongly
increases bending stiffness. Increased fillet radius reduces the assumed
stress concentration. The model does not calculate bolt bearing, bolt
preload, weld-group stress, weld-toe fatigue, local plate bending, contact, or
three-dimensional load sharing.

\subsection{Padeye}
\label{sec:padeye-definition}

The padeye represents a lifting point for a 500~kg working mass with a dynamic
factor of 3:

\begin{equation}
F_d=500(9.81)(3)=14\,715~\mathrm{N}.
\label{eq:padeye-load}
\end{equation}

The project-defined S275 properties are density 7850~kg/m$^3$, yield strength
275~MPa, Young's modulus 210~GPa, and Poisson ratio 0.30. The fixed pin-hole
diameter is 50~mm, and two gussets reinforce the lug. The criteria are
$FOS\geq2.5$ and $\delta\leq0.5$~mm.

\begin{table}[H]
\centering
\caption{Padeye design variables}
\label{tab:padeye-parameters}
\begin{tabularx}{\textwidth}{Xcrrr}
\toprule
\textbf{Variable} & \textbf{Unit} & \textbf{Lower} & \textbf{Baseline} &
\textbf{Upper}\\
\midrule
Base length & mm & 220 & 260 & 320\\
Base width & mm & 150 & 180 & 230\\
Base thickness & mm & 8 & 14 & 22\\
Lug height & mm & 140 & 180 & 230\\
Lug thickness & mm & 12 & 20 & 30\\
Neck width & mm & 80 & 100 & 140\\
Gusset height & mm & 70 & 110 & 120\\
Gusset thickness & mm & 6 & 10 & 18\\
Fillet radius & mm & 8 & 15 & 25\\
\bottomrule
\end{tabularx}
\end{table}

\subsubsection{Analytical formulation}
\label{sec:padeye-model}

The volume model combines the doubler plate, a tapered lug and circular crown,
the 50~mm hole subtraction, two triangular gussets, and a fillet
approximation. This representation follows the native generator closely enough
for a mass consistency check, but it is not an exact boundary-representation
calculation.

The structural model calculates four candidate stresses. The net-section
tension stress uses the remaining ligament at the hole and a
fillet-dependent stress factor. The nominal bearing stress is

\begin{equation}
\sigma_{\mathrm{bearing}}=
\frac{F_d}{d_p t_l},
\label{eq:padeye-bearing}
\end{equation}

where $d_p$ is pin diameter and $t_l$ is lug thickness. The shear-out
estimate uses the two load paths between the hole and the outer edge:

\begin{equation}
\tau_{\mathrm{out}}=
\frac{F_d}{2A_{\mathrm{shear}}}.
\label{eq:padeye-shearout}
\end{equation}

A lateral-bending estimate represents an assumed 25\% side-load component.
The reported padeye stress is

\begin{equation}
\sigma_p=\max\left(
\sigma_{\mathrm{net}},
\sigma_{\mathrm{bearing}},
\sqrt{3}\tau_{\mathrm{out}},
\sigma_{\mathrm{lateral}}
\right).
\label{eq:padeye-governing-stress}
\end{equation}

The displacement combines axial and lateral terms, and the factor of safety is
$FOS_p=\sigma_y/\sigma_p$. This model ranks candidates by several relevant
nominal modes. It does not model a separate pin, pin clearance, nonlinear
contact, local peak stress, weld geometry, or fatigue. These limitations
define the later FEA and design-code work.

\subsection{Stabilizer}
\label{sec:stabilizer-definition}

The stabilizer retains a NACA 0015 outer envelope with 800~mm span, 500~mm root
chord, 250~mm tip chord, 12-degree sweep, and 8~mm trailing edge. It has a
560~mm by 160~mm by 18~mm root flange and three ribs at 25\%, 50\%, and 75\%
span.

The project-defined Aluminum 6061-T6 properties match the bracket. The
prescribed resultant follows

\begin{equation}
F_d=\frac{1}{2}\rho_wV^2SC_L\gamma_d,
\label{eq:stabilizer-load-general}
\end{equation}

where seawater density is 1025~kg/m$^3$, speed is 8~m/s, planform area is
0.3~m$^2$, lift coefficient is 0.8, and the dynamic factor is 1.5. Thus

\begin{equation}
F_d=
\frac{1}{2}(1025)(8^2)(0.3)(0.8)(1.5)
=11\,808~\mathrm{N}.
\label{eq:stabilizer-load}
\end{equation}

Four pressure bands use relative weights 1.00, 0.93, 0.75, and 0.40. These are
declared engineering assumptions, not CFD data. The screening criteria are
$FOS\geq2.5$ and tip displacement no more than 5~mm.

\begin{table}[H]
\centering
\caption{Stabilizer design variables}
\label{tab:stabilizer-parameters}
\begin{tabularx}{\textwidth}{Xcrrr}
\toprule
\textbf{Variable} & \textbf{Unit} & \textbf{Lower} & \textbf{Baseline} &
\textbf{Upper}\\
\midrule
Skin thickness & mm & 4 & 7 & 12\\
Front spar position & \% chord & 25 & 30 & 38\\
Rear spar position & \% chord & 58 & 65 & 72\\
Front spar thickness & mm & 4 & 8 & 14\\
Rear spar thickness & mm & 4 & 6 & 12\\
Rib thickness & mm & 3 & 5 & 9\\
Root insert length & mm & 120 & 180 & 260\\
Root insert thickness & mm & 8 & 12 & 20\\
Root fillet radius & mm & 12 & 25 & 40\\
\bottomrule
\end{tabularx}
\end{table}

\subsubsection{Analytical formulation}
\label{sec:stabilizer-model}

The volume model combines the two-sided skin, front and rear spar webs, three
internal ribs, root insert, fixed flange, and root-fillet approximation. The
NACA thickness equation estimates the local spar height at each chord
position. The root section combines skin, spar, and insert contributions to
the second moment of area.

The root bending moment uses a resultant centroid at 40\% span:

\begin{equation}
M_r=0.40F_db.
\label{eq:stabilizer-moment}
\end{equation}

The model calculates a nominal root bending stress $\sigma_b$ and shear stress
$\tau$. A root-fillet factor $K_t$ modifies the bending term. The equivalent
stress is

\begin{equation}
\sigma_s=
\sqrt{(K_t\sigma_b)^2+3\tau^2}.
\label{eq:stabilizer-stress}
\end{equation}

The tip displacement estimate is

\begin{equation}
\delta_s=\frac{F_db^3}{8EI_{\mathrm{eff}}}.
\label{eq:stabilizer-displacement}
\end{equation}

The model does not derive the pressure field. It also omits torsional
deflection, shell buckling, panel crippling, local skin--spar stress, fatigue,
fluid--structure interaction, and root-bearing detail. It therefore supports
structural screening under a prescribed load, not hydrodynamic optimization.

\section{Design-Space Exploration and Optimization}
\label{sec:search-method}

\subsection{Evaluation and Failure Handling}
\label{sec:evaluation-contract}

Each evaluator returns a structured record with volume, mass, maximum stress,
factor of safety, maximum displacement, objective, status, and error. It first
checks that the part exists, all required parameters are present, values are
finite, values lie within bounds, and the requested evaluation mode is
available.

Invalid inputs and unavailable evaluators return a failed record with
objective $10^9$. The analytical optimizer can continue after such a record.
A CAD-only response has partial status because it can contain valid mass and
artifacts while structural metrics remain unavailable. The reserved
\code{open\_fea} mode deliberately returns a controlled failure. No Gmsh or
CalculiX result is fabricated.

\subsection{Latin-Hypercube Exploration}
\label{sec:sampling-method}

Each part used 60 Latin-hypercube samples with seed 42. The sampler operated in
the scaled unit cube. Each variable was divided into 60 intervals. A random
point was selected within each interval, and independently permuted variable
columns formed the multi-dimensional vectors. The inverse scaling in
Equation~\ref{eq:variable-scaling} produced physical values.

The sampler evaluated all points, retained failed records, and counted
feasible candidates. A feasible candidate met both the factor-of-safety and
displacement constraints. Promising feasible sampled designs supplied two of
the three optimization starts.

\subsection{Bounded Multi-Start Optimization}
\label{sec:optimization-method}

The pure-Python optimizer operated on scaled coordinates. Each part used three
starts:

\begin{enumerate}
    \item the configured baseline;
    \item a promising low-mass feasible sample; and
    \item a second promising feasible sample that provided another basin or
    parameter combination.
\end{enumerate}

Each run permitted at most 80 simplex iterations. The optimizer converted each
trial point to physical units, evaluated the part, cached repeated vectors,
and appended each unique evaluation to the run log. The simplex stopped at the
iteration limit or when the objective spread met the configured tolerance.

\begin{figure}[H]
\centering
\begin{tikzpicture}[node distance=1.05cm,every node/.style={font=\small}]
\node[startstop] (start) {Load part definition};
\node[process,below=of start] (baseline) {Evaluate baseline and\\60 seeded LHS points};
\node[process,below=of baseline] (starts) {Select baseline and\\two sampled starts};
\node[process,below=of starts] (nm) {Run bounded Nelder--Mead\\for each start};
\node[decision,below=1.35cm of nm] (valid) {Valid candidate?};
\node[process,below left=1.3cm and 2.6cm of valid] (fail) {Record failed row\\with $10^9$ objective};
\node[process,below right=1.3cm and 2.6cm of valid] (metrics) {Record metrics and\\penalized objective};
\node[process,below=3.0cm of valid] (filter) {Combine runs and filter\\explicitly for feasibility};
\node[startstop,below=of filter] (package) {Select baseline and finalists;\\create validation package};

\draw[arrow] (start) -- (baseline);
\draw[arrow] (baseline) -- (starts);
\draw[arrow] (starts) -- (nm);
\draw[arrow] (nm) -- (valid);
\draw[arrow] (valid) -- node[above left]{no} (fail);
\draw[arrow] (valid) -- node[above right]{yes} (metrics);
\draw[arrow] (fail) |- (filter);
\draw[arrow] (metrics) |- (filter);
\draw[arrow] (filter) -- (package);
\end{tikzpicture}
\caption{Sampling, optimization, and finalist-selection procedure}
\label{fig:optimization-flow}
\end{figure}

The selection step is deliberately separate from the optimizer's minimum
objective. The validation package sorts only feasible rows by mass. It retains
five finalists so that the lightest design does not become the only candidate
for later FEA.

\subsection{Results Recording and Review}
\label{sec:logging-method}

Each evaluated candidate was recorded with its part identifier, parameter
values, calculated metrics, feasibility status, objective, and any evaluation
error. Run-level metadata preserved the seed and search context. A dashboard
was used to inspect convergence, parameter activity, mass, constraint use, and
failed evaluations without changing the underlying results.

\section{CAD Generation and Boundary Definition}
\label{sec:cad-method}

\subsection{Versioned CAD Handoff}
\label{sec:fusion-bridge}

Each CAD request identifies the schema, component, run, iteration, evaluation
mode, parameter values, and required artifacts. A legacy form is retained for
the original bracket demonstration. Requests and responses are written
atomically, and each returned result is accepted only when its run and
iteration identifiers match the candidate under evaluation. This prevents a
partial or stale exchange from being interpreted as current evidence.

\begin{figure}[H]
\centering
\begin{tikzpicture}[
    node distance=1.4cm and 2.0cm,
    lifeline/.style={rectangle,draw=nemoblue,fill=nemolight,thick,
    minimum width=4.0cm,minimum height=0.9cm,align=center},
    msg/.style={-{Stealth[length=3mm]},thick,draw=nemoblue}
]
\node[lifeline] (external) {External Python process};
\node[lifeline,right=6.0cm of external] (fusion) {Fusion NEMOBridge};
\draw[dashed] (external.south) -- ++(0,-7.4);
\draw[dashed] (fusion.south) -- ++(0,-7.4);
\draw[msg] ($(external.south)+(0,-0.8)$) --
node[above,font=\small]{candidate request} ($(fusion.south)+(0,-0.8)$);
\draw[msg] ($(fusion.south)+(0,-2.0)$) --
node[above,font=\small]{validate schema and identifiers}
($(external.south)+(0,-2.0)$);
\draw[msg] ($(external.south)+(0,-3.2)$) --
node[above,font=\small]{parameters and artifact request}
($(fusion.south)+(0,-3.2)$);
\draw[msg] ($(fusion.south)+(0,-4.4)$) --
node[above,font=\small]{native CAD, mass, STEP, boundary tags}
($(external.south)+(0,-4.4)$);
\draw[msg] ($(fusion.south)+(0,-5.6)$) --
node[above,font=\small]{correlated response}
($(external.south)+(0,-5.6)$);
\node[process,below=6.2cm of external] (log) {Correlate response\\and log record};
\node[process,below=6.2cm of fusion] (partial) {Structural fields remain null\\until manual FEA};
\end{tikzpicture}
\caption{Serialized Fusion request and response sequence}
\label{fig:bridge-sequence}
\end{figure}

The handoff permits one active CAD request at a time. This serialized approach
favours traceability over throughput and prevents two candidates from
competing for the same exchange state.

\subsection{Native CAD Generation}
\label{sec:cad-generation}

The Fusion generator uses the same engineering definitions as the analytical
study while remaining within Fusion's separate Python environment.

\subsubsection{Bracket}

The bracket generator creates a rectangular baseplate, cuts four mounting
holes, creates two triangular ribs, joins the bodies, and applies four long
rib-root fillets. It checks that the result is one connected solid with
positive volume. It identifies the mounting-hole cylindrical faces as the
fixed-support role and the four upper sloping rib faces as the
equipment-load role.

\subsubsection{Padeye}

The padeye generator creates the doubler plate, tapered lug, circular crown,
pin bore, two root fillets, and two overlapping gussets. It requires one
connected solid. The underside of the doubler plate supplies the
fixed-support role. The cylindrical pin-hole region supplies the
pin-bearing role.

\subsubsection{Stabilizer}

The stabilizer generator creates the root flange, lofted NACA outer envelope,
root fillet, hollow interior, front and rear spars, three ribs, root insert,
and pressure-band splits. It requires one connected solid. The root flange
provides the fixed support. Four surface regions provide the pressure bands,
and a tip region provides the displacement monitor.

Autodesk exposes physical-property and export functions through the public API
\cite{autodeskPhysicalProperties,autodeskExport}. \NEMO\ records mass and
volume from the generated component and exports STEP for downstream exchange.
STL is optional because it is a tessellated artifact rather than the preferred
solid exchange format.

\subsection{Semantic Boundary Definition}
\label{sec:boundary-method}

For each required role, the generator stores one or more face signatures.
Available fields include area, centroid, normal, bounding box, and cylindrical
geometry. The generator fails the sweep acceptance check if a required role
has no matching face.

This metadata has three functions. First, it documents what the generator
believed the load or support region to be. Second, it helps the engineer find
the same region in the manual study. Third, it makes topology problems visible
during a sweep.

It does not automatically apply the FEA boundary condition. The engineer must
still display the faces, compare them with the intended load path, and reject
an incorrect selection.

\section{Verification and Validation}
\label{sec:verification-validation-method}

\subsection{CAD Generator Verification}
\label{sec:cad-sweep-method}

Each generator was tested with its configured baseline and 20 seeded
Latin-hypercube vectors. The acceptance criteria for each of the 21 cases were:

\begin{enumerate}
    \item the request completed without an uncontrolled exception;
    \item the result contained exactly one connected solid;
    \item the physical volume was positive;
    \item STEP export completed;
    \item boundary metadata export completed; and
    \item every required semantic boundary role matched at least one face.
\end{enumerate}

The sweep does not solve FEA. It tests whether the generator and boundary
contract remain usable across the bounded design space. It is intentionally
separate from the offline test suite because it requires an open Fusion
session and mutates the shared live channel.

\subsection{Finalist Validation Preparation}
\label{sec:validation-package-method}

For each component, the baseline and five light feasible candidates were
assembled into a validation set. Each record retained its parameters,
analytical predictions, intended material and boundaries, CAD status and mass,
and fields for later FEA inputs, outputs, convergence evidence, and engineering
disposition. This preserved the connection between a solver image and the
candidate to which it belonged.

\subsection{Bracket Static Stress Studies}
\label{sec:bracket-poc-method}

The bracket proof of concept continued from the analytical finalist set into
Fusion Static Stress. Aluminum 6061 was assigned to the generated bracket, the
four mounting-hole surfaces represented the initial fixed-support
idealization, and the four upper sloping rib faces represented the equipment
load interface. The intended total downward load was 1472~N, obtained by
rounding the analytical design load of 1471.5~N to the integer input used for
the study.

Two exploratory studies were reviewed. Their material assignment, boundary
regions, mesh information, result magnitudes, and result locations were
compared with the engineering definition. The review also checked whether the
applied force was a total load rather than a load repeated on every selected
face. The images and findings from both studies are presented and interpreted
in Section~\ref{sec:bracket-fea-result}.

\subsection{Static Stress Assessment Criteria}
\label{sec:fea-quality-procedure}

Study credibility was assessed from the material properties, applied-load
magnitude and direction, support and contact idealizations, mesh order and
refinement, reaction balance, result locations, and stability under mesh
change. Maximum von Mises stress, minimum factor of safety, and maximum
displacement were considered only after these input and numerical checks.
Peaks at constrained edges were treated cautiously because they can represent
local singular behaviour. Even an acceptable linear-static study would cover
only the stated load case and would not establish fatigue, buckling, weld,
bolt, corrosion, impact, or classification compliance.

\subsection{Extension to Padeye FEA}
\label{sec:padeye-fea-method}

The padeye study will follow the same evidence sequence with part-specific
boundaries and result regions. The underside of the doubler plate will provide
the first fixed-support idealization. The 14\,715~N design load will act
through the pin-bearing region.

An initial linear screening model can use a distributed bearing load on the
hole surface. A later detailed model should include a separate pin, clearance,
contact, and any required nonlinear behavior. The review regions are the
loaded hole, crown, net section, shear-out ligaments, lug root, gusset roots,
doubler interface, and weld lines. Reaction balance must close to
14\,715~N in the applied direction.

The final design assessment must also check the selected lifting code, load
angle range, dynamic factors, material certification, weld class, fabrication
tolerance, inspection, proof load, fatigue, and supporting structure.

\subsection{Extension to Stabilizer FEA}
\label{sec:stabilizer-fea-method}

The stabilizer study will fix the root-flange interface. The 11\,808~N
resultant will be distributed over four pressure bands. The band forces are
calculated by normalizing the weights 1.00, 0.93, 0.75, and 0.40. Pressure on
each band equals its force divided by its actual CAD face area.

The study will monitor tip displacement and inspect the root fillet, front and
rear spars, skin--spar intersections, rib junctions, insert termination, root
flange, and trailing edge. Shell or solid element choice must represent the
thin skins and local joints adequately. Buckling and fatigue require separate
studies.

The prescribed load distribution is an interim structural envelope. A later
hydrodynamic stage must define operational conditions and derive pressure from
an accepted method. The structural model should then include positive and
negative load cases, torsion, actuator loads, and relevant vessel motions.

\subsection{Reproducibility and Method Limitations}
\label{sec:method-verification}

Software verification covered engineering-definition loading, bounds,
baseline feasibility, qualitative model sensitivity, sampling, optimizer
behaviour, controlled failures, results recording, CAD handoff, and finalist
preparation. Fusion sweeps were conducted separately because they depend on an
active desktop application; they checked geometry and boundary evidence but
did not solve FEA.

Exact finalist parameters, analytical results, CAD masses, and solver-study
evidence were retained so that the tables and figures in this report could be
traced to the corresponding candidate. The method improves reproducibility,
but it does not remove uncertainty in the analytical assumptions, candidate
selection, CAD boundary identification, or manual solver setup.

\screeningwarning

The method was designed to reduce candidate count and improve traceability. It
was not designed to certify marine hardware. A future engineering decision
must use verified loads, applicable codes, detailed FEA, uncertainty
assessment, fatigue and fabrication checks, physical testing where required,
and qualified approval.
